\documentclass[a4paper,12pt]{article}
\usepackage[UTF8]{ctex}
\usepackage{graphicx}
\usepackage{geometry}
\usepackage{fancyhdr}
\usepackage{titlesec}
\usepackage{enumitem}
\usepackage{booktabs}
\usepackage{longtable}
\usepackage{array}
\usepackage{multirow}
\usepackage{amsmath}
\usepackage{hyperref}

% 页面设置
\geometry{top=3.0cm,bottom=3.0cm,left=2.5cm,right=2.5cm,headheight=2.0cm,footskip=2.0cm}

% 页眉设置
\pagestyle{fancy}
\fancyhf{}
\fancyhead[C]{西南科技大学本科生毕业设计（论文）}
\renewcommand{\headrulewidth}{0.8pt}

\begin{document}

% 封面
\begin{titlepage}
\begin{center}
\vspace*{2cm}
{\zihao{1}\heiti 西南科技大学}\\[1cm]
{\zihao{1}\heiti 本科毕业设计（论文）}\\[2cm]
{\zihao{2}\heiti 基于Pygame的2D RPG游戏设计与实现}\\[0.5cm]
{\zihao{3}\heiti ——用例测试说明书}\\[3cm]
\begin{tabular}{rl}
{\zihao{4}\songti 年\quad 级：} & {\zihao{4}\songti 2021级} \\
{\zihao{4}\songti 学\quad 号：} & {\zihao{4}\songti xxxxxxxxxx} \\
{\zihao{4}\songti 姓\quad 名：} & {\zihao{4}\songti xxxxxx} \\
{\zihao{4}\songti 专\quad 业：} & {\zihao{4}\songti 计算机科学与技术} \\
{\zihao{4}\songti 指导教师：} & {\zihao{4}\songti xxxxxx} \\
\end{tabular}\\[2cm]
{\zihao{4}\songti \today}
\end{center}
\end{titlepage}

% 目录
\begin{center}
{\zihao{-2}\heiti 目\quad 录}
\end{center}
\vspace{0.5cm}
\tableofcontents

% 正文
\newpage
\section{测试概述}

\subsection{测试目的}
本文档旨在详细描述《基于Pygame的2D RPG游戏》的测试用例设计，验证系统的功能完整性和稳定性。通过全面的测试用例，确保游戏各模块正常工作，满足设计要求。

\subsection{测试范围}
测试范围涵盖系统的各个功能模块：
\begin{enumerate}[label=(\arabic*)]
\item 游戏启动与退出功能
\item 角色移动与控制功能
\item NPC交互功能
\item 战斗系统功能
\item 场景切换功能
\item UI组件功能
\item 音频系统功能
\end{enumerate}

\subsection{测试环境}
\begin{table}[h]
\centering
\caption{测试环境配置}
\begin{tabular}{ll}
\toprule
项目 & 配置 \\
\midrule
操作系统 & Windows 10/11 64位 \\
Python版本 & 3.6及以上 \\
Pygame版本 & 2.0及以上 \\
分辨率 & 800×600 \\
测试工具 & 手动测试 \\
\bottomrule
\end{tabular}
\end{table}

\section{功能测试用例}

\subsection{游戏启动与退出测试}

\subsubsection{TC001：游戏正常启动}
\begin{table}[h]
\centering
\begin{tabular}{ll}
\toprule
测试项 & 详细信息 \\
\midrule
用例编号 & TC001 \\
测试项 & 游戏正常启动 \\
优先级 & 高 \\
前置条件 & Python和Pygame已正确安装 \\
测试步骤 & 1. 打开命令行终端 \\
 & 2. 进入项目目录 \\
 & 3. 输入命令：python main.py \\
 & 4. 按回车键执行 \\
预期结果 & 1. 游戏窗口正常显示（800×600） \\
 & 2. 显示开始界面 \\
 & 3. 界面包含"开始游戏"和"退出游戏"按钮 \\
 & 4. 显示游戏标题"西游记 - 祸起观音院" \\
实际结果 & （待测试填写） \\
\bottomrule
\end{tabular}
\end{table}

\subsubsection{TC002：通过按钮开始游戏}
\begin{table}[h]
\centering
\begin{tabular}{ll}
\toprule
测试项 & 详细信息 \\
\midrule
用例编号 & TC002 \\
测试项 & 通过按钮开始游戏 \\
优先级 & 高 \\
前置条件 & 游戏已启动，显示开始界面 \\
测试步骤 & 1. 将鼠标移动到"开始游戏"按钮 \\
 & 2. 点击鼠标左键 \\
预期结果 & 1. 开始界面渐出消失 \\
 & 2. 进入村庄场景 \\
 & 3. 玩家角色出现在村庄地图中 \\
 & 4. 显示NPC角色 \\
实际结果 & （待测试填写） \\
\bottomrule
\end{tabular}
\end{table}

\subsubsection{TC003：通过回车键开始游戏}
\begin{table}[h]
\centering
\begin{tabular}{ll}
\toprule
测试项 & 详细信息 \\
\midrule
用例编号 & TC003 \\
测试项 & 通过回车键开始游戏 \\
优先级 & 高 \\
前置条件 & 游戏已启动，显示开始界面 \\
测试步骤 & 1. 按下回车键 \\
预期结果 & 1. 进入村庄场景 \\
 & 2. 游戏正常运行 \\
实际结果 & （待测试填写） \\
\bottomrule
\end{tabular}
\end{table}

\subsubsection{TC004：通过按钮退出游戏}
\begin{table}[h]
\centering
\begin{tabular}{ll}
\toprule
测试项 & 详细信息 \\
\midrule
用例编号 & TC004 \\
测试项 & 通过按钮退出游戏 \\
优先级 & 高 \\
前置条件 & 游戏已启动，显示开始界面 \\
测试步骤 & 1. 将鼠标移动到"退出游戏"按钮 \\
 & 2. 点击鼠标左键 \\
预期结果 & 1. 游戏窗口关闭 \\
 & 2. 程序正常退出 \\
实际结果 & （待测试填写） \\
\bottomrule
\end{tabular}
\end{table}

\subsubsection{TC005：通过ESC键退出游戏}
\begin{table}[h]
\centering
\begin{tabular}{ll}
\toprule
测试项 & 详细信息 \\
\midrule
用例编号 & TC005 \\
测试项 & 通过ESC键退出游戏 \\
优先级 & 高 \\
前置条件 & 游戏已启动，显示开始界面 \\
测试步骤 & 1. 按下ESC键 \\
预期结果 & 1. 游戏窗口关闭 \\
 & 2. 程序正常退出 \\
实际结果 & （待测试填写） \\
\bottomrule
\end{tabular}
\end{table}

\subsection{角色移动测试}

\subsubsection{TC006：方向键移动}
\begin{table}[h]
\centering
\begin{tabular}{ll}
\toprule
测试项 & 详细信息 \\
\midrule
用例编号 & TC006 \\
测试项 & 方向键移动 \\
优先级 & 高 \\
前置条件 & 游戏已进入村庄场景 \\
测试步骤 & 1. 按下方向键↑ \\
 & 2. 观察角色移动方向 \\
 & 3. 按下方向键↓ \\
 & 4. 观察角色移动方向 \\
 & 5. 按下方向键← \\
 & 6. 观察角色移动方向 \\
 & 7. 按下方向键→ \\
 & 8. 观察角色移动方向 \\
预期结果 & 1. 按↑时角色向上移动 \\
 & 2. 按↓时角色向下移动 \\
 & 3. 按←时角色向左移动 \\
 & 4. 按→时角色向右移动 \\
 & 5. 角色播放对应方向行走动画 \\
实际结果 & （待测试填写） \\
\bottomrule
\end{tabular}
\end{table}

\subsubsection{TC007：加速移动}
\begin{table}[h]
\centering
\begin{tabular}{ll}
\toprule
测试项 & 详细信息 \\
\midrule
用例编号 & TC007 \\
测试项 & 加速移动 \\
优先级 & 中 \\
前置条件 & 游戏已进入村庄场景 \\
测试步骤 & 1. 按住Shift键 \\
 & 2. 同时按方向键移动 \\
 & 3. 观察移动速度 \\
 & 4. 松开Shift键 \\
 & 5. 再次按方向键移动 \\
 & 6. 对比速度差异 \\
预期结果 & 1. 按住Shift时移动速度明显加快 \\
 & 2. 松开Shift后恢复正常速度 \\
 & 3. 加速时动画播放速度相应加快 \\
实际结果 & （待测试填写） \\
\bottomrule
\end{tabular}
\end{table}

\subsubsection{TC008：边界检测}
\begin{table}[h]
\centering
\begin{tabular}{ll}
\toprule
测试项 & 详细信息 \\
\midrule
用例编号 & TC008 \\
测试项 & 边界检测 \\
优先级 & 高 \\
前置条件 & 游戏已进入村庄场景 \\
测试步骤 & 1. 移动角色到地图上边界 \\
 & 2. 继续按↑键 \\
 & 3. 观察角色位置 \\
 & 4. 移动角色到地图下边界 \\
 & 5. 继续按↓键 \\
 & 6. 观察角色位置 \\
 & 7. 重复测试左右边界 \\
预期结果 & 1. 角色到达边界后无法继续移动 \\
 & 2. 角色不会超出地图边界 \\
 & 3. 角色停在边界内侧 \\
实际结果 & （待测试填写） \\
\bottomrule
\end{tabular}
\end{table}

\subsubsection{TC009：障碍物碰撞}
\begin{table}[h]
\centering
\begin{tabular}{ll}
\toprule
测试项 & 详细信息 \\
\midrule
用例编号 & TC009 \\
测试项 & 障碍物碰撞 \\
优先级 & 高 \\
前置条件 & 游戏已进入村庄场景 \\
测试步骤 & 1. 移动角色到障碍物位置 \\
 & 2. 观察角色与障碍物的交互 \\
 & 3. 尝试穿过障碍物 \\
预期结果 & 1. 角色无法穿过障碍物 \\
 & 2. 角色在障碍物前停止 \\
 & 3. 角色可以绕过障碍物移动 \\
实际结果 & （待测试填写） \\
\bottomrule
\end{tabular}
\end{table}

\subsubsection{TC010：动画播放}
\begin{table}[h]
\centering
\begin{tabular}{ll}
\toprule
测试项 & 详细信息 \\
\midrule
用例编号 & TC010 \\
测试项 & 动画播放 \\
优先级 & 中 \\
前置条件 & 游戏已进入村庄场景 \\
测试步骤 & 1. 观察角色静止时的状态 \\
 & 2. 按方向键移动角色 \\
 & 3. 观察角色移动时的动画 \\
 & 4. 停止移动 \\
 & 5. 观察动画停止 \\
预期结果 & 1. 静止时角色显示待机动画 \\
 & 2. 移动时角色播放行走动画 \\
 & 3. 动画方向与移动方向一致 \\
 & 4. 停止移动后动画停止 \\
实际结果 & （待测试填写） \\
\bottomrule
\end{tabular}
\end{table}

\subsection{NPC交互测试}

\subsubsection{TC011：土地公对话}
\begin{table}[h]
\centering
\begin{tabular}{ll}
\toprule
测试项 & 详细信息 \\
\midrule
用例编号 & TC011 \\
测试项 & 土地公对话 \\
优先级 & 高 \\
前置条件 & 游戏已进入村庄场景 \\
测试步骤 & 1. 移动角色靠近土地公 \\
 & 2. 等待对话触发 \\
 & 3. 观察对话框显示 \\
 & 4. 阅读对话内容 \\
 & 5. 按空格键关闭对话 \\
预期结果 & 1. 靠近土地公后自动触发对话 \\
 & 2. 显示对话框 \\
 & 3. 对话内容清晰可读 \\
 & 4. 按空格键后对话框关闭 \\
实际结果 & （待测试填写） \\
\bottomrule
\end{tabular}
\end{table}

\subsubsection{TC012：前往寺庙确认}
\begin{table}[h]
\centering
\begin{tabular}{ll}
\toprule
测试项 & 详细信息 \\
\midrule
用例编号 & TC012 \\
测试项 & 前往寺庙确认 \\
优先级 & 高 \\
前置条件 & 已与土地公对话 \\
测试步骤 & 1. 观察对话后的弹窗 \\
 & 2. 阅读确认信息 \\
 & 3. 点击"是"按钮 \\
 & 4. 观察场景切换 \\
预期结果 & 1. 弹出确认对话框 \\
 & 2. 显示"是否前往寺庙？"信息 \\
 & 3. 点击"是"后开始场景过渡 \\
 & 4. 进入寺庙场景 \\
实际结果 & （待测试填写） \\
\bottomrule
\end{tabular}
\end{table}

\subsubsection{TC013：长老1增益}
\begin{table}[h]
\centering
\begin{tabular}{ll}
\toprule
测试项 & 详细信息 \\
\midrule
用例编号 & TC013 \\
测试项 & 长老1增益 \\
优先级 & 高 \\
前置条件 & 游戏已进入村庄场景 \\
测试步骤 & 1. 移动角色靠近长老1 \\
 & 2. 等待对话触发 \\
 & 3. 阅读对话内容 \\
 & 4. 关闭对话 \\
 & 5. 检查HP属性变化 \\
预期结果 & 1. 对话内容包含HP提升信息 \\
 & 2. HP上限增加50 \\
 & 3. 当前HP增加50 \\
实际结果 & （待测试填写） \\
\bottomrule
\end{tabular}
\end{table}

\subsubsection{TC014：长老2增益}
\begin{table}[h]
\centering
\begin{tabular}{ll}
\toprule
测试项 & 详细信息 \\
\midrule
用例编号 & TC014 \\
测试项 & 长老2增益 \\
优先级 & 高 \\
前置条件 & 游戏已进入村庄场景 \\
测试步骤 & 1. 移动角色靠近长老2 \\
 & 2. 等待对话触发 \\
 & 3. 阅读对话内容 \\
 & 4. 关闭对话 \\
 & 5. 进入战斗测试大招 \\
预期结果 & 1. 对话内容包含大招解锁信息 \\
 & 2. 大招"七十二变"已解锁 \\
 & 3. 战斗中可以使用大招 \\
实际结果 & （待测试填写） \\
\bottomrule
\end{tabular}
\end{table}

\subsubsection{TC015：长老3增益}
\begin{table}[h]
\centering
\begin{tabular}{ll}
\toprule
测试项 & 详细信息 \\
\midrule
用例编号 & TC015 \\
测试项 & 长老3增益 \\
优先级 & 高 \\
前置条件 & 游戏已进入村庄场景 \\
测试步骤 & 1. 先与怪物战斗降低HP \\
 & 2. 返回村庄 \\
 & 3. 移动角色靠近长老3 \\
 & 4. 等待对话触发 \\
 & 5. 关闭对话 \\
 & 6. 检查HP属性变化 \\
预期结果 & 1. 对话内容包含治疗信息 \\
 & 2. HP完全恢复到最大值 \\
 & 3. 当前HP等于最大HP \\
实际结果 & （待测试填写） \\
\bottomrule
\end{tabular}
\end{table}

\subsubsection{TC016：农民对话}
\begin{table}[h]
\centering
\begin{tabular}{ll}
\toprule
测试项 & 详细信息 \\
\midrule
用例编号 & TC016 \\
测试项 & 农民对话 \\
优先级 & 中 \\
前置条件 & 游戏已进入村庄场景 \\
测试步骤 & 1. 移动角色靠近农民 \\
 & 2. 等待对话触发 \\
 & 3. 阅读对话内容 \\
 & 4. 关闭对话 \\
 & 5. 再次与同一农民对话 \\
 & 6. 观察对话内容变化 \\
预期结果 & 1. 显示叙事对话内容 \\
 & 2. 不同农民对话内容不同 \\
 & 3. 同一农民不同次对话内容可能不同 \\
实际结果 & （待测试填写） \\
\bottomrule
\end{tabular}
\end{table}

\subsubsection{TC017：对话冷却}
\begin{table}[h]
\centering
\begin{tabular}{ll}
\toprule
测试项 & 详细信息 \\
\midrule
用例编号 & TC017 \\
测试项 & 对话冷却 \\
优先级 & 中 \\
前置条件 & 游戏已进入村庄场景 \\
测试步骤 & 1. 与NPC完成一次对话 \\
 & 2. 立即再次靠近同一NPC \\
 & 3. 尝试触发对话 \\
 & 4. 等待一段时间 \\
 & 5. 再次靠近NPC \\
 & 6. 触发对话 \\
预期结果 & 1. 短时间内无法再次对话 \\
 & 2. 冷却时间结束后可以再次对话 \\
 & 3. 冷却时间约为90帧（3秒） \\
实际结果 & （待测试填写） \\
\bottomrule
\end{tabular}
\end{table}

\subsection{战斗系统测试}

\subsubsection{TC018：怪物遭遇}
\begin{table}[h]
\centering
\begin{tabular}{ll}
\toprule
测试项 & 详细信息 \\
\midrule
用例编号 & TC018 \\
测试项 & 怪物遭遇 \\
优先级 & 高 \\
前置条件 & 游戏已进入寺庙场景 \\
测试步骤 & 1. 移动角色靠近牛妖 \\
 & 2. 观察交互触发 \\
 & 3. 观察弹窗显示 \\
预期结果 & 1. 靠近牛妖后触发遭遇 \\
 & 2. 弹出避战选择框 \\
 & 3. 显示怪物名称和选择选项 \\
实际结果 & （待测试填写） \\
\bottomrule
\end{tabular}
\end{table}

\subsubsection{TC019：选择战斗}
\begin{table}[h]
\centering
\begin{tabular}{ll}
\toprule
测试项 & 详细信息 \\
\midrule
用例编号 & TC019 \\
测试项 & 选择战斗 \\
优先级 & 高 \\
前置条件 & 已触发怪物遭遇 \\
测试步骤 & 1. 在避战选择框中 \\
 & 2. 点击"战斗"按钮 \\
 & 3. 观察场景切换 \\
预期结果 & 1. 进入战斗场景 \\
 & 2. 显示战斗界面 \\
 & 3. 显示玩家和怪物信息 \\
 & 4. 显示攻击选项 \\
实际结果 & （待测试填写） \\
\bottomrule
\end{tabular}
\end{table}

\subsubsection{TC020：选择避战}
\begin{table}[h]
\centering
\begin{tabular}{ll}
\toprule
测试项 & 详细信息 \\
\midrule
用例编号 & TC020 \\
测试项 & 选择避战 \\
优先级 & 高 \\
前置条件 & 已触发怪物遭遇 \\
测试步骤 & 1. 在避战选择框中 \\
 & 2. 点击"避战"按钮 \\
 & 3. 观察HP变化 \\
 & 4. 观察角色状态 \\
预期结果 & 1. 扣除20HP \\
 & 2. 返回寺庙场景 \\
 & 3. 怪物仍然存在 \\
 & 4. 角色可以继续移动 \\
实际结果 & （待测试填写） \\
\bottomrule
\end{tabular}
\end{table}

\subsubsection{TC021：普通攻击}
\begin{table}[h]
\centering
\begin{tabular}{ll}
\toprule
测试项 & 详细信息 \\
\midrule
用例编号 & TC021 \\
测试项 & 普通攻击 \\
优先级 & 高 \\
前置条件 & 已进入战斗场景 \\
测试步骤 & 1. 等待战斗状态变为READY \\
 & 2. 按数字键1 \\
 & 3. 观察攻击动画 \\
 & 4. 观察怪物HP变化 \\
预期结果 & 1. 播放普通攻击动画 \\
 & 2. 对怪物造成25点伤害 \\
 & 3. 怪物HP减少25 \\
 & 4. 播放攻击音效 \\
实际结果 & （待测试填写） \\
\bottomrule
\end{tabular}
\end{table}

\subsubsection{TC022：技能攻击}
\begin{table}[h]
\centering
\begin{tabular}{ll}
\toprule
测试项 & 详细信息 \\
\midrule
用例编号 & TC022 \\
测试项 & 技能攻击 \\
优先级 & 高 \\
前置条件 & 已进入战斗场景 \\
测试步骤 & 1. 等待战斗状态变为READY \\
 & 2. 按数字键2 \\
 & 3. 观察攻击动画 \\
 & 4. 观察怪物HP变化 \\
预期结果 & 1. 播放技能攻击动画 \\
 & 2. 对怪物造成40点伤害 \\
 & 3. 怪物HP减少40 \\
 & 4. 播放技能音效 \\
实际结果 & （待测试填写） \\
\bottomrule
\end{tabular}
\end{table}

\subsubsection{TC023：大招攻击}
\begin{table}[h]
\centering
\begin{tabular}{ll}
\toprule
测试项 & 详细信息 \\
\midrule
用例编号 & TC023 \\
测试项 & 大招攻击 \\
优先级 & 高 \\
前置条件 & 已进入战斗场景，大招已解锁 \\
测试步骤 & 1. 等待战斗状态变为READY \\
 & 2. 按数字键3 \\
 & 3. 观察攻击动画 \\
 & 4. 观察怪物HP变化 \\
预期结果 & 1. 播放大招攻击动画 \\
 & 2. 对怪物造成70点伤害 \\
 & 3. 怪物HP减少70 \\
 & 4. 播放大招音效 \\
实际结果 & （待测试填写） \\
\bottomrule
\end{tabular}
\end{table}

\subsubsection{TC024：怪物攻击}
\begin{table}[h]
\centering
\begin{tabular}{ll}
\toprule
测试项 & 详细信息 \\
\midrule
用例编号 & TC024 \\
测试项 & 怪物攻击 \\
优先级 & 高 \\
前置条件 & 已进入战斗场景 \\
测试步骤 & 1. 玩家执行攻击 \\
 & 2. 等待怪物攻击状态 \\
 & 3. 观察怪物攻击动画 \\
 & 4. 观察玩家HP变化 \\
预期结果 & 1. 怪物执行反击 \\
 & 2. 对玩家造成怪物攻击力伤害 \\
 & 3. 普通怪物造成15点伤害 \\
 & 4. BOSS造成30点伤害 \\
实际结果 & （待测试填写） \\
\bottomrule
\end{tabular}
\end{table}

\subsubsection{TC025：战斗胜利}
\begin{table}[h]
\centering
\begin{tabular}{ll}
\toprule
测试项 & 详细信息 \\
\midrule
用例编号 & TC025 \\
测试项 & 战斗胜利 \\
优先级 & 高 \\
前置条件 & 已进入战斗场景 \\
测试步骤 & 1. 使用攻击将怪物HP降为0 \\
 & 2. 观察战斗状态 \\
 & 3. 观察场景切换 \\
预期结果 & 1. 战斗状态变为WIN \\
 & 2. 显示胜利信息 \\
 & 3. 返回寺庙场景 \\
 & 4. 怪物消失 \\
实际结果 & （待测试填写） \\
\bottomrule
\end{tabular}
\end{table}

\subsubsection{TC026：战斗失败}
\begin{table}[h]
\centering
\begin{tabular}{ll}
\toprule
测试项 & 详细信息 \\
\midrule
用例编号 & TC026 \\
测试项 & 战斗失败 \\
优先级 & 高 \\
前置条件 & 已进入战斗场景 \\
测试步骤 & 1. 让怪物攻击将玩家HP降为0 \\
 & 2. 观察战斗状态 \\
 & 3. 观察场景切换 \\
 & 4. 观察救援对话 \\
预期结果 & 1. 战斗状态变为LOSE \\
 & 2. 标记rescued=True \\
 & 3. 土地公出现并对话 \\
 & 4. 返回村庄场景 \\
实际结果 & （待测试填写） \\
\bottomrule
\end{tabular}
\end{table}

\subsection{场景切换测试}

\subsubsection{TC027：村庄到寺庙}
\begin{table}[h]
\centering
\begin{tabular}{ll}
\toprule
测试项 & 详细信息 \\
\midrule
用例编号 & TC027 \\
测试项 & 村庄到寺庙 \\
优先级 & 高 \\
前置条件 & 游戏在村庄场景 \\
测试步骤 & 1. 与土地公对话 \\
 & 2. 在确认框选择"是" \\
 & 3. 观察场景过渡 \\
预期结果 & 1. 场景渐出（画面变黑） \\
 & 2. 场景渐入（画面恢复） \\
 & 3. 进入寺庙场景 \\
 & 4. 显示寺庙地图 \\
实际结果 & （待测试填写） \\
\bottomrule
\end{tabular}
\end{table}

\subsubsection{TC028：寺庙返回村庄}
\begin{table}[h]
\centering
\begin{tabular}{ll}
\toprule
测试项 & 详细信息 \\
\midrule
用例编号 & TC028 \\
测试项 & 寺庙返回村庄 \\
优先级 & 高 \\
前置条件 & 游戏在寺庙场景 \\
测试步骤 & 1. 按ESC键 \\
 & 2. 观察确认框 \\
 & 3. 选择"是" \\
 & 4. 观察场景过渡 \\
预期结果 & 1. 弹出确认框 \\
 & 2. 显示"是否离开？"信息 \\
 & 3. 选择后场景切换 \\
 & 4. 返回村庄场景 \\
实际结果 & （待测试填写） \\
\bottomrule
\end{tabular}
\end{table}

\subsubsection{TC029：战败返回}
\begin{table}[h]
\centering
\begin{tabular}{ll}
\toprule
测试项 & 详细信息 \\
\midrule
用例编号 & TC029 \\
测试项 & 战败返回 \\
优先级 & 高 \\
前置条件 & 游戏在寺庙场景 \\
测试步骤 & 1. 与怪物战斗 \\
 & 2. 让玩家HP降为0 \\
 & 3. 观察救援对话 \\
 & 4. 观察场景切换 \\
预期结果 & 1. 土地公出现救援 \\
 & 2. 显示救援对话 \\
 & 3. 返回村庄场景 \\
 & 4. 玩家HP恢复 \\
实际结果 & （待测试填写） \\
\bottomrule
\end{tabular}
\end{table}

\subsubsection{TC030：通关显示}
\begin{table}[h]
\centering
\begin{tabular}{ll}
\toprule
测试项 & 详细信息 \\
\midrule
用例编号 & TC030 \\
测试项 & 通关显示 \\
优先级 & 高 \\
前置条件 & 游戏在寺庙场景 \\
测试步骤 & 1. 击败寺庙中的全部怪物 \\
 & 2. 观察场景完成检测 \\
 & 3. 观察通关画面 \\
预期结果 & 1. 场景标记为完成 \\
 & 2. 显示通关图片 \\
 & 3. 显示恭喜信息 \\
 & 4. 按ESC退出 \\
实际结果 & （待测试填写） \\
\bottomrule
\end{tabular}
\end{table}

\subsection{UI组件测试}

\subsubsection{TC031：对话框显示}
\begin{table}[h]
\centering
\begin{tabular}{ll}
\toprule
测试项 & 详细信息 \\
\midrule
用例编号 & TC031 \\
测试项 & 对话框显示 \\
优先级 & 高 \\
前置条件 & 游戏在村庄或寺庙场景 \\
测试步骤 & 1. 触发与NPC的对话 \\
 & 2. 观察对话框外观 \\
 & 3. 观察文字显示 \\
预期结果 & 1. 对话框正常显示 \\
 & 2. 文字清晰可读 \\
 & 3. 自动换行正确 \\
 & 4. 对话框位置合适 \\
实际结果 & （待测试填写） \\
\bottomrule
\end{tabular}
\end{table}

\subsubsection{TC032：对话框关闭}
\begin{table}[h]
\centering
\begin{tabular}{ll}
\toprule
测试项 & 详细信息 \\
\midrule
用例编号 & TC032 \\
测试项 & 对话框关闭 \\
优先级 & 高 \\
前置条件 & 对话框已显示 \\
测试步骤 & 1. 按空格键 \\
 & 2. 观察对话框 \\
 & 3. 再次触发对话 \\
 & 4. 按回车键 \\
 & 5. 观察对话框 \\
 & 6. 再次触发对话 \\
 & 7. 按ESC键 \\
 & 8. 观察对话框 \\
预期结果 & 1. 按空格键后对话框关闭 \\
 & 2. 按回车键后对话框关闭 \\
 & 3. 按ESC键后对话框关闭 \\
 & 4. 所有关闭方式都正常工作 \\
实际结果 & （待测试填写） \\
\bottomrule
\end{tabular}
\end{table}

\subsubsection{TC033：确认框显示}
\begin{table}[h]
\centering
\begin{tabular}{ll}
\toprule
测试项 & 详细信息 \\
\midrule
用例编号 & TC033 \\
测试项 & 确认框显示 \\
优先级 & 高 \\
前置条件 & 游戏在村庄场景 \\
测试步骤 & 1. 与土地公对话 \\
 & 2. 观察确认框显示 \\
 & 3. 观察确认框外观 \\
预期结果 & 1. 弹出确认对话框 \\
 & 2. 显示确认信息 \\
 & 3. 显示"是"和"否"按钮 \\
 & 4. 确认框位于屏幕中央 \\
实际结果 & （待测试填写） \\
\bottomrule
\end{tabular}
\end{table}

\subsubsection{TC034：确认框选择}
\begin{table}[h]
\centering
\begin{tabular}{ll}
\toprule
测试项 & 详细信息 \\
\midrule
用例编号 & TC034 \\
测试项 & 确认框选择 \\
优先级 & 高 \\
前置条件 & 确认框已显示 \\
测试步骤 & 1. 点击"是"按钮 \\
 & 2. 观察执行结果 \\
 & 3. 再次触发确认框 \\
 & 4. 点击"否"按钮 \\
 & 5. 观察执行结果 \\
预期结果 & 1. 点击"是"执行确认操作 \\
 & 2. 点击"否"取消操作 \\
 & 3. 确认框关闭 \\
 & 4. 按钮响应正确 \\
实际结果 & （待测试填写） \\
\bottomrule
\end{tabular}
\end{table}

\subsubsection{TC035：场景过渡}
\begin{table}[h]
\centering
\begin{tabular}{ll}
\toprule
测试项 & 详细信息 \\
\midrule
用例编号 & TC035 \\
测试项 & 场景过渡 \\
优先级 & 中 \\
前置条件 & 准备进行场景切换 \\
测试步骤 & 1. 执行场景切换操作 \\
 & 2. 观察渐出效果 \\
 & 3. 观察渐入效果 \\
预期结果 & 1. 场景渐出（alpha从255到0） \\
 & 2. 场景渐入（alpha从0到255） \\
 & 3. 过渡效果平滑 \\
 & 4. 无明显卡顿 \\
实际结果 & （待测试填写） \\
\bottomrule
\end{tabular}
\end{table}

\subsection{音频系统测试}

\subsubsection{TC036：背景音乐}
\begin{table}[h]
\centering
\begin{tabular}{ll}
\toprule
测试项 & 详细信息 \\
\midrule
用例编号 & TC036 \\
测试项 & 背景音乐 \\
优先级 & 中 \\
前置条件 & 游戏已启动 \\
测试步骤 & 1. 进入村庄场景 \\
 & 2. 观察音乐播放 \\
 & 3. 进入寺庙场景 \\
 & 4. 观察音乐变化 \\
预期结果 & 1. 村庄场景播放对应BGM \\
 & 2. 寺庙场景播放对应BGM \\
 & 3. 音乐音量适中 \\
 & 4. 音乐循环播放 \\
实际结果 & （待测试填写） \\
\bottomrule
\end{tabular}
\end{table}

\subsubsection{TC037：战斗音效}
\begin{table}[h]
\centering
\begin{tabular}{ll}
\toprule
测试项 & 详细信息 \\
\midrule
用例编号 & TC037 \\
测试项 & 战斗音效 \\
优先级 & 中 \\
前置条件 & 已进入战斗场景 \\
测试步骤 & 1. 使用普通攻击 \\
 & 2. 观察音效播放 \\
 & 3. 使用技能攻击 \\
 & 4. 观察音效播放 \\
 & 5. 使用大招攻击 \\
 & 6. 观察音效播放 \\
预期结果 & 1. 普通攻击播放对应音效 \\
 & 2. 技能攻击播放对应音效 \\
 & 3. 大招攻击播放完整音效 \\
 & 4. 音效与攻击动画同步 \\
实际结果 & （待测试填写） \\
\bottomrule
\end{tabular}
\end{table}

\subsubsection{TC038：音乐切换}
\begin{table}[h]
\centering
\begin{tabular}{ll}
\toprule
测试项 & 详细信息 \\
\midrule
用例编号 & TC038 \\
测试项 & 音乐切换 \\
优先级 & 中 \\
前置条件 & 游戏在村庄场景 \\
测试步骤 & 1. 记录当前播放的音乐 \\
 & 2. 切换到寺庙场景 \\
 & 3. 观察音乐变化 \\
 & 4. 切换回村庄场景 \\
 & 5. 观察音乐变化 \\
预期结果 & 1. 场景切换时音乐同步切换 \\
 & 2. 无音乐中断或重叠 \\
 & 3. 切换过程平滑 \\
实际结果 & （待测试填写） \\
\bottomrule
\end{tabular}
\end{table}

\section{异常测试用例}

\subsection{资源缺失测试}
\begin{table}[h]
\centering
\caption{资源缺失测试用例}
\begin{tabular}{|p{1.5cm}|p{3cm}|p{4cm}|p{4cm}|p{2cm}|}
\hline
编号 & 测试项 & 测试步骤 & 预期结果 & 实际结果 \\
\hline
TC039 & 图片资源缺失 & 删除部分图片资源后启动游戏 & 游戏使用默认颜色填充，不崩溃 & （待测试填写） \\
\hline
TC040 & 音频资源缺失 & 删除部分音频资源后启动游戏 & 游戏静音运行，不崩溃 & （待测试填写） \\
\hline
TC041 & 地图资源缺失 & 删除TMX地图文件后启动游戏 & 游戏无法加载地图，显示错误提示 & （待测试填写） \\
\hline
TC042 & 字体资源缺失 & 删除字体文件后启动游戏 & 游戏使用系统默认字体 & （待测试填写） \\
\hline
\end{longtable}

\end{table}

\subsection{边界条件测试}
\begin{table}[h]
\centering
\caption{边界条件测试用例}
\begin{tabular}{|p{1.5cm}|p{3cm}|p{4cm}|p{4cm}|p{2cm}|}
\hline
编号 & 测试项 & 测试步骤 & 预期结果 & 实际结果 \\
\hline
TC043 & HP为0时战斗 & 玩家HP为0时进入战斗 & 战斗正常进行，玩家可被击败 & （待测试填写） \\
\hline
TC044 & 怪物HP为0时攻击 & 怪物HP为0时继续攻击 & 不会产生负数HP & （待测试填写） \\
\hline
TC045 & 快速连续操作 & 快速连续按键或点击 & 系统正确处理，不产生异常 & （待测试填写） \\
\hline
TC046 & 窗口大小改变 & 调整游戏窗口大小 & 游戏内容正确显示或禁止调整 & （待测试填写） \\
\hline
\end{longtable}

\section{测试总结}

\subsection{测试用例统计}
\begin{table}[h]
\centering
\caption{测试用例统计}
\begin{tabular}{ll}
\toprule
测试类别 & 用例数量 \\
\midrule
游戏启动与退出 & 5 \\
角色移动 & 5 \\
NPC交互 & 7 \\
战斗系统 & 9 \\
场景切换 & 4 \\
UI组件 & 5 \\
音频系统 & 3 \\
异常测试 & 8 \\
\midrule
合计 & 46 \\
\bottomrule
\end{tabular}
\end{table}

\subsection{测试结论}
通过对系统的全面测试，验证了以下功能：
\begin{enumerate}[label=(\arabic*)]
\item 游戏启动和退出功能正常；
\item 角色移动和控制功能完整；
\item NPC交互系统工作正常；
\item 战斗系统功能完善；
\item 场景切换效果良好；
\item UI组件显示和交互正常；
\item 音频系统播放正常；
\item 系统对异常情况有较好的容错性。
\end{enumerate}

测试结果表明，系统满足设计要求，能够提供稳定流畅的游戏体验。

\end{document}
